It is to be regretted that, during the three years when the theological professors of the United Church were at Augsburg, Prof. Sverdrup did not teach Dogmatics—and, we may add, Symbolics. In Dogmatics one would have received the best that Sverdrup had to offer, together with a scholarly grounding of the many doctrinal propositions that both he and the leading conservative-liberal Lutheran theologians in Europe accepted, but of which our local dogmatic brood here partly knew little, and partly did not wish to know anything. In Symbolics one would have received a scholarly account of the ecclesiastical position he occupied.

If Sverdrup had taught Dogmatics and Symbolics in the years 1890–93, which he did both before and afterward, the historical oppositions that came to light in the days of strife might perhaps have looked somewhat different. As matters stood, the opportunity for an academic scholarly discussion of systematic theology and for a historical treatment of the church’s confessional symbols was cut off. And Sverdrup was too much a man of scholarship to find the whole of Lutheran Dogmatics in Job and Daniel. He kept to “the Text.”

Owing to the unhappy church conflict in Norway, the word “scholarly” has acquired a bad sound in certain of our church circles. But just as there is much that passes itself off as Christianity which is not Christianity, so there is much that is called scholarship which is nothing of the kind. We have used the word in its proper sense. As secretary of the theological faculty at Augsburg in 1875, Sverdrup calls the faculty “the bearer of scholarship.” “We strive to make ourselves organs of a reciprocal interaction between the congregation and scholarship, which we regard as a fundamental condition for an ecclesiastical development, and which only the free church gives occasion for in full measure.” (Conf. Annual Report 1875, 63f.)

It is especially in systematic theology that this reciprocal interaction will bear ripe fruit. But the one who teaches this subject ought to know what congregation is, and what scholarship is. Sverdrup knew the one as well as the other; but in the faculty there were those who taught what they had inherited from St. Louis and from Mt. Airy: the seventeenth century’s elementary conception of congregation, and the eighteenth century’s conception of scholarship, which had to lead to confusions.

Sverdrup has been called a disciple of Gisle Johnson. But surely not with justice. For the fact that a man hears a professor for some years does not of itself make him that professor’s disciple. Sverdrup studied as long abroad as he did in Christiania, and his writings, so far as theology is concerned, bear witness far more to German influence than to Norwegian. Nor did he owe anything to Fr. Petersen, a man for whose eclectic theology Sverdrup had no esteem. The judgment that Bishop Bang passes on Petersen in his Reminiscences Sverdrup would no doubt have been able to sign as well. Neither did Sverdrup owe anything to Krogh-Tonning. Sverdrup’s students know well enough what sort of rough handling Krogh-Tonning repeatedly received in the old theological classroom at Augsburg. When others found Tonning’s doctrine of the church in order—indeed, so good that it had to be crammed—Sverdrup repeatedly let his students understand that Tonning was ready for departure to the Catholic Church, where he too at last came to harbor.

In Symbolics Sverdrup was the same broad-minded man as in Dogmatics. For him the point was to find what was historically warranted in the different branches of the church, not to go on crusade after “deviating doctrines.” It is characteristic that he branded chapters 4–13 in Graul’s Symbolics as “trifling pedantries,” and referred the students, for a correspondingly concise but factual treatment of the sects, to Prof. Gumlich’s (d. 1902, Berlin) Symbolics.

Nor can it be said that Sverdrup was a disciple of Caspari. He tended rather toward the psychological direction represented by Delitzsch and inherited by Buhl.

If we are to connect Sverdrup’s name with any theological school, then our choice must fall upon the Erlangen school. In his view of history, Scripture, the Confession, theological method, and the church, Sverdrup was essentially an Erlangen theologian. So far from standing isolated in his ecclesiastical-theological position, he had the most excellent predecessors to support him: eminent theologians such as Harless, Höfling, Hofmann, Schöberlein, Thomasius, Delitzsch, Luthardt, Theodosius Harnack, Frank, Plitt, Beischwitz. Among his contemporaries, in whom the Erlangen school entered a new period of splendor, may be named an authority in the history of the Old Covenant such as Köhler, a patristic scholar and New Testament exegete such as Zahn, church historians such as Hauck and Kolde, dogmaticians such as Seeberg and Ihmels.

It is characteristic of our circumstances that an Erlangen scion such as Seeberg’s well-known book, At the Professor’s Chair into the Twentieth Century, received a very laudatory review from Prof. Sverdrup (in Folkebladet), while it became the object of wholly different treatment in Lutheraneren, where a union man such as Pastor Ole Nilsen and an anti-missionary such as Pastor Gjære occupied 40–50 columns in order to show that Prof. Seeberg was a false teacher—despite the fact that Norwegian orthodox theologians such as Prof. Odland and Provost Gustav Jensen had also given the book the best recommendations.

Erlangen theology came to Augsburg in the seventies. The attempt that was made in the nineties to exchange it for the seventeenth century’s repristination theology failed.

Whether, or to what degree, Sverdrup worked creatively in theology shall not be decided here. We may concur with Bishop Bang when he says: “Over there [among us] one has not been able to cultivate theological scholarship for scholarship’s own sake.” And it is at any rate not too much to say when he adds: “One has indeed shown theological scholarship a serious labor in practical-ecclesiastical interest, and this has deposited fruits in the literature that are worthy of all honor.”

In the Collected Writings the weight has been laid upon the practical. This will become evident from what follows.

(Volume I. Sketches and Essays)

The first volume is by now so well known, at least by prospectus, that we need not communicate the titles of the various essays or treat the separate contents. About 300 pages in this volume are reprinted from Kvartal-Skrift, which is now very difficult to obtain except at great cost. The reprinting is something many will welcome.

As typical of Sverdrup’s treatment of the Old Testament narratives, where the point was edifying exposition, the piece “The Prophet Samuel” deserves special attention. Eli is presented as a priest utterly devoid of insight into spiritual things. Everything he did was wrong. The pious praying Hannah he accused of being drunk, since he did not know the difference between the effects of wine and the effects of the Spirit. His house was undisciplined. The announcement of punishment from the Lord he listened to in an apparently pious resignation which in reality was helpless indifferentism. Prophets had to be raised up; the worship at the temple had been made worldly.

The ethical essays on the sanctification of Sunday and on the common school also deserve attention, essays occasioned by the Synod’s radical understanding of the day of rest and by its disparaging judgment upon the common school because it does not teach religion. Here too Sverdrup shows himself, as we knew him, a good American citizen, who defends a democratic institution—and an evangelical Christian, who upholds the free nature of religion. No doubt the German Missourians have their own elementary schools, but these schools stand far behind our common schools. Experience will show that Sverdrup’s contribution in the matter has not been put to shame. The essay on Sunday ought to be read by the many among us who share the Reformed view of the Sabbath.

Of a practical-dogmatic kind, and important for understanding why Sverdrup had so little use for the Synod’s “pure doctrine,” is the third section of the book. The essay on “Justification by Faith and the Justification of the Teachers” is as instructive and as necessary for us now as it was thirty years ago. That Sverdrup in these essays spoke a timely word will appear from what Prof. Madsen at the University of Copenhagen writes (in Nielsen’s Church Lexicon for the North) about the Bornholmers, and from Prof. Fr. Nielsen’s article on the same subject in Herzog-Hauck’s Protestant Realencyclopedia. From these articles one will also understand why the justification of the world is still preached in various places among us. Those especially who preach this misleading doctrine are those who have inherited from Missouri, or who read Rosenius without criticism, unaware that he fell back into the doctrine that it is only unbelief that condemns.

The third section is of a church-historical character. It contains, among other things, four essays still little known. One of them, “Criticism in History,” has until now existed only in manuscript. Prof. Helland states that, to judge from the handwriting, it belongs to the very earliest years after Sverdrup’s arrival in America. The title as well as the content is characteristic. Those who think that Sverdrup’s speech about criticism began in the nineties will here, if not elsewhere, see that they are mistaken.

Also rather unknown in our circles is the essay “The Norwegian Lutheran Church in America,” which Sverdrup had written in French for the Lutheran church paper Le Temoignage, and which from there, in translation, was taken up in Vestlandsposten (1897).

Of distinctive interest is the sketch “The Lutheran Free Church.” It was written in 1907 at the request of the Census Bureau in Washington, D.C. I join myself to Prof. Helland’s words: “It will have an increased interest from the fact that it is Prof. Sverdrup’s last word on what, according to his understanding, ‘criticism’ is and is to be. In that respect this little sketch may be said to form his testament.”

What we are most grateful for is “Memories from Norway,” a fragment of a manuscript composed during a stay in Norway in the summer of 1886. Here Sverdrup sketches a whole series of Norwegian churchmen from the side of revival. We set great value upon this, since Sverdrup seldom wrote characterizations of more recent men or told anything about them in class.

I cannot understand why the review in Lutheraneren says that “Memories from Norway” “is very unfinished, more nearly a rough draft, and could therefore best have been omitted from a book such as this.” In Norway it is precisely this piece to which attention has especially been drawn as interesting. Thus Luthersk Kirketidende, Vor Kirke og Kultur (Jens Gleditsch), Morgenbladet (Bishop Bang).

Let me be allowed to conclude these remarks on the first volume by citing a couple of statements from Norway. Pastor Lars Dahle says in Norsk Missionstidende: “Everything (in the first volume) bears the stamp of the author’s powerful, earnest, and logical personality.” Bishop Bang says: “On every page one will find more vivid and strong traces of the author’s warm personality and strong-willed energy.” Of the whole undertaking Luthersk Kirketidende says: “The Norwegian church people ought to support this. Prof. Sverdrup was both so considerable a churchman and theologian that one ought to make oneself acquainted with his work.” Bishop Bang recommends the work in the best terms.
